\secnum{5}
\section{Low Pass Filter}\label{sec:lpf}
Bộ lọc thông thấp cho qua tín hiệu từ $0\,\text{Hz}$ (DC) đến tần số cắt trên $f_H$ và chặn tín hiệu có tần số cao hơn. Tại $f_H$ độ lợi giảm $3\,\text{dB}$ so với độ lợi băng thông (biên độ còn $1/\sqrt{2} \approx 70.7\%$).

\given{fig5_08}{0.5\linewidth}{Low pass filter principles}

\subsection{Inverting low pass filter}\label{sec:lpf-a}
Mạch Hình~\ref{fig:5-09}: $R_1 = 1\,\text{k}\Omega$, $R_2 = 10\,\text{k}\Omega$ song song $C = 1\,\text{nF}$ ở nhánh hồi tiếp, nguồn \texttt{VAC} $2\,\text{V}_{AC}$ ($0\,\text{V}_{DC}$), ngõ $(+)$ nối đất.

\given{fig5_09}{0.5\linewidth}{Inverting low pass filter\label{fig:5-09}}
\dsn{ex3.5.dsn}

\subsubsection*{Tính toán lý thuyết}
\begin{theorybox}
\textbf{Hàm truyền.} Trở kháng nhánh hồi tiếp $Z_f = R_2 \parallel \dfrac{1}{j\omega C} = \dfrac{R_2}{1 + j\omega R_2 C}$. Mạch đảo có $H = -Z_f/R_1$:
\[
H(j\omega) = \frac{V_{out}}{V_{in}} = -\frac{R_2}{R_1}\cdot\frac{1}{1 + j\omega R_2 C}
= -\frac{A_0}{1 + j\,f/f_H},
\qquad
|H| = \frac{A_0}{\sqrt{1 + (f/f_H)^2}}.
\]
\textbf{Độ lợi tại $0\,\text{Hz}$} (tụ hở mạch):
\[
A_0 = \frac{R_2}{R_1} = \frac{10\,\text{k}}{1\,\text{k}} = 10
\;\Rightarrow\; 20\log_{10}10 = 20\,\text{dB},\qquad V_{out} = 10 \times 2 = 20\,\text{V}.
\]
\textbf{Tần số cắt:}
\[
f_H = \frac{1}{2\pi R_2 C} = \frac{1}{2\pi \times 10^{4} \times 10^{-9}} \approx 15.92\,\text{kHz}.
\]
\textbf{Độ lợi tại $f_H$:}
\[
|H(f_H)| = \frac{10}{\sqrt{2}} \approx 7.07
\;\Rightarrow\; 20\log_{10}7.07 \approx 17.0\,\text{dB} = 20 - 3\,\text{dB},\qquad V_{out} \approx 14.14\,\text{V}.
\]
\textbf{Tại $50\,\text{kHz}$} (cuối dải quét): $|H| = 10/\sqrt{1 + (50/15.92)^2} \approx 3.03$, $V_{out} \approx 6.07\,\text{V}$.
\end{theorybox}

\subsubsection*{Mô phỏng PSpice}
\simcfg{AC Sweep}{Linear, Start $=1\,\text{Hz}$, Stop $=50\,\text{kHz}$, 200 points}
Đặt Voltage Marker ở ngõ ra; trong Probe dùng \textit{Cursor} tìm điểm $V_{out} = 14.14\,\text{V}$ (hoặc vẽ \texttt{DB(V(out)/V(in))} và tìm điểm $17\,\text{dB}$) để đọc $f_H$.
\simresult
\simfig{ex35_sch}{Sơ đồ mạch inverting low pass filter (PSpice)}
\simfig{ex35_ac}{Đồ thị AC Sweep, đánh dấu tần số cắt $f_H$ (cursor)}

\begin{table}[H]
\centering
\begin{tabular}{l c c}
\toprule
\textbf{Đại lượng} & \textbf{Lý thuyết} & \textbf{Mô phỏng}\\
\midrule
$V_{out}$ tại $\approx 0\,\text{Hz}$ (V) & 20     & \fillcell\\
Gain tại $0\,\text{Hz}$ (V/V)            & 10     & \fillcell\\
Gain tại $0\,\text{Hz}$ (dB)             & 20     & \fillcell\\
$f_H$ (kHz)                              & 15.92  & \fillcell\\
$V_{out}$ tại $f_H$ (V)                  & 14.14  & \fillcell\\
Gain tại $f_H$ (V/V)                     & 7.07   & \fillcell\\
Gain tại $f_H$ (dB)                      & 17.0   & \fillcell\\
\bottomrule
\end{tabular}
\caption{Inverting low pass filter ($V_{IN}=2\,\text{V}$)}
\end{table}
% TODO: kiểm lại sau khi có số liệu mô phỏng
\begin{remarkbox}
\begin{itemize}[nosep]
  \item Ở tần số thấp tụ $C$ gần như hở mạch, mạch là bộ khuếch đại đảo độ lợi $R_2/R_1 = 10$ ($V_{out} = 20\,\text{V}$). Tần số tăng thì trở kháng tụ $1/(2\pi f C)$ giảm, nối tắt dần $R_2$, độ lợi giảm.
  \item Tần số cắt đọc trên đồ thị (điểm độ lợi giảm $3\,\text{dB}$, $V_{out} \approx 14.14\,\text{V}$) xấp xỉ $16\,\text{kHz}$, khớp với $f_H = 1/(2\pi R_2 C)$.
  \item Sau $f_H$ độ lợi giảm với độ dốc $-20\,\text{dB/decade}$ (bộ lọc bậc 1). Vì dải quét tuyến tính chỉ đến $50\,\text{kHz}$ ($\approx 3f_H$) nên đồ thị chưa thấy rõ độ dốc này; quét kiểu Decade sẽ thấy rõ hơn.
  \item $R_1$ chỉ quyết định độ lợi, $f_H$ chỉ phụ thuộc $R_2C$: có thể chỉnh độ lợi và tần số cắt tương đối độc lập.
\end{itemize}
\end{remarkbox}

\subsection{Non-inverting low pass filter \bonus}\label{sec:lpf-b}
Tính $R$, $C$ để độ lợi băng thông bằng $10$ và tần số cắt bằng mục~\ref{sec:lpf-a} (Hình~\ref{fig:5-12}: điện trở $1\,\text{k}\Omega$ nối tiếp ngõ $(+)$, $R_{11} = 1\,\text{k}\Omega$ từ ngõ $(-)$ xuống đất, $R \parallel C$ ở nhánh hồi tiếp).

\given{fig5_12}{0.45\linewidth}{Non-inverting low pass filter\label{fig:5-12}}
\dsn{ex3.5b.dsn}

\subsubsection*{Tính toán thiết kế}
\begin{theorybox}
\textbf{Hàm truyền.} Không có dòng vào ngõ $(+)$ nên điện trở $1\,\text{k}\Omega$ nối tiếp không gây sụt áp: $V_{(+)} = V_{in}$. Với $Z_f = R \parallel \dfrac{1}{j\omega C}$:
\[
H(j\omega) = 1 + \frac{Z_f}{R_{11}} = 1 + \frac{R/R_{11}}{1 + j\omega RC}.
\]
\textbf{Độ lợi băng thông ($f \to 0$):}
\[
A_0 = 1 + \frac{R}{R_{11}} = 10 \;\Rightarrow\; R = 9R_{11} = 9\,\text{k}\Omega.
\]
\textbf{Tần số cắt.} Đặt $x = \omega RC$, với $R/R_{11} = 9$: $H = \dfrac{10 + jx}{1 + jx}$, $|H|^2 = \dfrac{100 + x^2}{1 + x^2}$. Điểm $-3\,\text{dB}$: $|H|^2 = A_0^2/2 = 50$:
\[
100 + x^2 = 50 + 50x^2 \;\Rightarrow\; x = \sqrt{50/49} \approx 1.010
\;\Rightarrow\; f_c = \frac{1.010}{2\pi RC}.
\]
Cho $f_c = f_H = 15.92\,\text{kHz}$ (bằng mục trước):
\[
C = \frac{1.010}{2\pi R f_c} = \frac{1.010}{2\pi \times 9000 \times 15\,915} \approx 1.12\,\text{nF}.
\]
(Công thức gần đúng $f_c \approx 1/(2\pi RC)$ cho $C \approx 1.11\,\text{nF}$, sai khác dưới $1\%$.) Chọn $C = 1.12\,\text{nF}$, kiểm tra lại: $f_c = 1.010/(2\pi \times 9\,\text{k} \times 1.12\,\text{n}) \approx 15.95\,\text{kHz}$.

\textbf{Khác biệt so với mạch đảo:} khi $f \to \infty$, $Z_f \to 0$ nên $|H| \to 1$ ($0\,\text{dB}$), không về $0$. Tại $50\,\text{kHz}$: $|H| \approx 3.16$ ($V_{out} \approx 6.3\,\text{V}$).
\end{theorybox}
\begin{table}[H]
\centering
\begin{tabular}{l c}
\toprule
\textbf{Linh kiện} & \textbf{Giá trị chọn}\\
\midrule
Điện trở nối tiếp ngõ $(+)$  & $1\,\text{k}\Omega$\\
$R_{11}$ (ngõ $(-)$ xuống đất) & $1\,\text{k}\Omega$\\
$R$ (hồi tiếp)               & $9\,\text{k}\Omega$\\
$C$ (hồi tiếp)               & $1.12\,\text{nF}$\\
\bottomrule
\end{tabular}
\caption{Giá trị thiết kế non-inverting low pass filter}
\end{table}
\begin{dapsobox}
$R = 9\,\text{k}\Omega$, \quad $C \approx 1.12\,\text{nF}$ \quad $\Rightarrow$ \quad $A_0 = 10$ ($20\,\text{dB}$), \quad $f_c \approx 15.9\,\text{kHz}$.
\end{dapsobox}

\subsubsection*{Mô phỏng PSpice}
\simcfg{AC Sweep}{Linear, Start $=1\,\text{Hz}$, Stop $=50\,\text{kHz}$, 200 points; nguồn \texttt{VAC} $2\,\text{V}$}
\simresult
\simfig{ex35b_sch}{Sơ đồ mạch non-inverting low pass filter (PSpice)}
\simfig{ex35b_ac}{Đồ thị AC Sweep, đánh dấu tần số cắt (cursor)}
\begin{table}[H]
\centering
\begin{tabular}{l c c}
\toprule
\textbf{Đại lượng} & \textbf{Thiết kế} & \textbf{Mô phỏng}\\
\midrule
$V_{out}$ tại $\approx 0\,\text{Hz}$ (V)         & 20    & \fillcell\\
Gain băng thông (V/V)                            & 10    & \fillcell\\
Tần số cắt $-3\,\text{dB}$ (kHz)                 & 15.95 & \fillcell\\
$V_{out}$ tại $f_c$ (V)                          & 14.14 & \fillcell\\
$V_{out}$ tại $50\,\text{kHz}$ (V)               & 6.3   & \fillcell\\
\bottomrule
\end{tabular}
\caption{Non-inverting low pass filter}
\end{table}
% TODO: kiểm lại sau khi có số liệu mô phỏng
\begin{remarkbox}
\begin{itemize}[nosep]
  \item Với $R = 9\,\text{k}\Omega$, $C = 1.12\,\text{nF}$, mạch đạt độ lợi băng thông $10$ và tần số cắt $\approx 16\,\text{kHz}$, giống mạch đảo; khác biệt là tín hiệu ra cùng pha với tín hiệu vào ở tần số thấp.
  \item Ở tần số cao độ lợi tiến về $1$ chứ không về $0$ (vì mạch không đảo luôn có độ lợi $\ge 1$), nên khả năng chặn tần số cao kém hơn mạch đảo: tại $50\,\text{kHz}$ ngõ ra $\approx 6.3\,\text{V}$ so với $\approx 6.07\,\text{V}$ ở mạch đảo, và chênh lệch càng lớn khi tần số tăng.
  \item Muốn bộ lọc không đảo suy giảm về $0$, thường đặt mạch RC thụ động ở ngõ vào $(+)$ rồi dùng op-amp khuếch đại không đảo.
\end{itemize}
\end{remarkbox}
